\documentclass[10pt,letterpaper]{article}
\usepackage[margin=0.85in]{geometry}
\usepackage{amsmath,amssymb,amsthm,mathtools,bm,booktabs,longtable,enumitem}
\usepackage[T1]{fontenc}
\usepackage{lmodern,microtype}
\usepackage[colorlinks=true,linkcolor=blue,urlcolor=blue,citecolor=blue]{hyperref}
\newtheorem{assumption}{Assumption}
\newtheorem{proposition}{Proposition}
\newtheorem{theorem}{Theorem}
\newtheorem{corollary}{Corollary}
\newtheorem{remark}{Remark}
\newcommand{\R}{\mathbb R}
\newcommand{\B}{\mathbb B}
\newcommand{\E}{\mathcal E}
\newcommand{\Z}{\mathcal Z}
\newcommand{\I}{\mathcal I}
\newcommand{\U}{\mathcal U}
\newcommand{\W}{\mathcal W}
\newcommand{\norm}[1]{\left\lVert#1\right\rVert}
\newcommand{\supp}[2]{\sigma_{#1}(#2)}
\newcommand{\Reach}{\operatorname{Reach}}
\newcommand{\dist}{\operatorname{dist}}
\newcommand{\diag}{\operatorname{diag}}
\newcommand{\He}{\operatorname{He}}
\newcommand{\cl}{\operatorname{cl}}
\title{From a Multistage Vehicle Observer to an Output-Feedback\protect\\
Robust Predictive Barrier and One Nominal CLF}
\author{Pan Dynamics Collision Avoidance Research}
\date{October 3, 2026; terminal set revised October 4, 2026;\protect\\
implemented tube, handling envelope and CLF row revised October 5, 2026}
\begin{document}
\maketitle
\begin{abstract}
We derive an information-state predictive safety controller from the existing
covariant NRMM observer certificate. Its links are explicit: a positive
comparison system bounds estimation error; an intersection update defines the
current possible true states; a causal output-feedback prediction propagates
these sets; support functions and interval reserves certify rectangle clearance;
a terminal set that only requires the ego to separate from the target and to
stay on the road bounds the clearance after the horizon when the relative
velocity stays constant; and the same path-following CLF provides
disturbance-dependent recovery. Each sample solves its own safety problem; no
recursive feasibility is claimed. One linearization is compatible with the
per-sample safety statement if its remainder bounds hold over the entire
decision and uncertainty domain. We also derive a vanishing proximal weight that
preserves CLF dissipation in one convex optimization. These are mathematical
design results, with finite numerical checks, not a claim that the present
certainty-equivalent controller already implements or certifies them.
\end{abstract}

\section{Claim, scope, and the single-controller design}
Let $t_k=kh$ and let $u_k$ be held on $[t_k,t_{k+1})$. We seek
\begin{equation}\label{eq:goal}
\dist\bigl(\mathcal V_E(t),\mathcal V_C(t)\bigr)\ge d_0,
\qquad t\ge t_0,
\end{equation}
for every admitted initial error, sensor error, plant disturbance and target
motion. Here $\mathcal V_E,\mathcal V_C$ are the full vehicle rectangles and
$d_0=0.20\,\mathrm m$ is the present design clearance, not an uncertainty bound.
The nominal behavior is the given path at its desired cruise speed, with free
longitudinal phase. There is one controller and one nominal CLF throughout the
encounter and recovery. The terminal set is a condition on the last predicted
node, not a second command mode.

The general use of online observer bounds in robust output-feedback MPC is
established prior work \cite{Kohler}. The present connection specializes that
foundation to the Sharma/NRMM vehicle model, its covariant cascade, relative
rectangle geometry, a predictive slack budget, and the existing one-CLF
architecture. The generic observer--tube--MPC chain itself is not a novelty claim.
Sharma et al.\ \cite{Sharma} supply the original vehicle/observer framework;
the repository uses a modified cascade, whose certificate must be verified on
its own equations. Huang et al.\ \cite{Huang} supply the predictive-slack
construction. Their exact optimal-value results are not silently transferred to
a retained suboptimal budget or an approximate nonlinear program.

\begin{assumption}[Declared model and information contract]\label{ass:contract}
The well-posed continuous plant and sensor maps, the sampled observer, and their
disturbance sets are specified. Initial errors have finite certified bounds.
Truth remains in the declared model domain, and each observer, geometry or
integration certificate is used only on its valid domain. Sensor timing, maximum admitted outages,
noise and hold errors satisfy their declared bounds. The target has constant
tangential acceleration $A_C$ and sideslip $\beta_C$, or an explicitly enclosed
model defect. No unknown future measurement is replaced by its nominal value.
\end{assumption}
For deterministic guarantees, Gaussian noise simulations alone are insufficient:
an untruncated Gaussian has no finite all-realization amplitude bound. A
probabilistic noise contract leads to a probabilistic safety claim with an
explicit time-uniform failure budget, not \eqref{eq:goal} for all realizations.
The positive-speed NRMM inverse is not valid through a stop. Constant negative
acceleration cannot be extrapolated through zero speed under this contract.
Likewise, finite-domain observer constants do not establish an infinite-time
claim after truth leaves that domain. The theorem below includes these premises;
they must be enforced or replaced by a valid extended model.

\section{Observer stability produces current state sets}
\subsection{Positive comparison, including the transient}
In the existing cascade let $e_v$ be the body-velocity error and
$e_T=[e_\rho^\top,e_q^\top,e_s^\top]^\top$. For an observer frequency $\omega_o>0$
and $P_o\succ0$, set
\begin{equation}
D_o=\diag(1,\omega_o^{-1},\omega_o^{-2}),\quad
Q_o=(D_oP_oD_o)\otimes I_2,\quad W_T=\sqrt{e_T^\top Q_o e_T}.
\end{equation}
The existing continuous-core derivation has the form
\begin{equation}\label{eq:comparison}
D^+\chi\le H\chi+d(t),\qquad
\chi=\begin{bmatrix}\norm{e_v}\\ W_T\end{bmatrix},\quad
H=\begin{bmatrix}-k_v&0\\c_{Tv}&-\lambda_T\end{bmatrix},
\quad k_v,\lambda_T>0,\ c_{Tv}\ge0,\ d(t)\ge0.
\end{equation}
All constants include the declared true-domain and disturbance terms. Define
\begin{equation}\label{eq:envelope}
b(t)=e^{H(t-t_0)}b_0+
\int_{t_0}^{t}e^{H(t-s)}\bar d(s)\,ds,
\qquad b_0\ge\chi(t_0),\quad \bar d\ge d.
\end{equation}
\begin{proposition}[Current error containment]\label{prop:observer}
Under \eqref{eq:comparison}, $\chi(t)\le b(t)$ componentwise. Consequently
\begin{equation}\label{eq:observerSet}
e_v\in b_v(t)\B_2,\qquad e_T\in\E_T(t):=
\{e:e^\top Q_o e\le b_T(t)^2\}.
\end{equation}
\end{proposition}
\begin{proof}
$H$ is Metzler, so $e^{H\tau}$ is entrywise nonnegative for $\tau\ge0$.
The comparison principle, followed by variation of constants, gives
\eqref{eq:envelope}. The two sets are the definitions of the bounded norms.
\end{proof}
For constant $\bar d$, $b(t)=e^{H\tau}b_0+H^{-1}(e^{H\tau}-I)\bar d$.
The lower-left entry of $e^{H\tau}$ is
\begin{equation}
c_{Tv}\frac{e^{-k_v\tau}-e^{-\lambda_T\tau}}{\lambda_T-k_v},
\quad\text{or }c_{Tv}\tau e^{-k_v\tau}\text{ if }\lambda_T=k_v.
\end{equation}
The steady-state bound $-H^{-1}\bar d$ cannot replace the transient unless it
also contains the initial error. The copositive weights
$w=-H^{-\top}{\bf1}>0$ give $w^\top H=-{\bf1}^\top$, providing an independent
certificate for the comparison system.

For any vector $a$, the exact ellipsoid support is
\begin{equation}\label{eq:support}
\supp{\E_T(t)}{a}=b_T(t)\sqrt{a^\top Q_o^{-1}a}.
\end{equation}
In particular the Euclidean component radii are
\begin{equation}\label{eq:components}
b_\rho=\sqrt{(P_o^{-1})_{11}}b_T,\quad
b_q=\omega_o\sqrt{(P_o^{-1})_{22}}b_T,\quad
b_s=\omega_o^2\sqrt{(P_o^{-1})_{33}}b_T.
\end{equation}
These follow by maximizing each directional projection; the maximizing error
is $b_TQ_o^{-1}a/\sqrt{a^\top Q_o^{-1}a}$. Keeping the joint ellipsoid preserves
correlation that separate component boxes discard.

\subsection{Sampled realization and the remaining ego channels}
The production observer is sampled and numerically integrated. A continuous
stability theorem alone is insufficient for its timestamped error. Use its
validated finite-time enclosure, or prove a sampled comparison update of the form
\begin{equation}\label{eq:sampled}
b_{k+1}\le \bar b_{1|k}:=M_kb_k+d_k^{\rm obs},\qquad M_k\ge0,
\end{equation}
where $d_k^{\rm obs}$ encloses sensor noise, held-signal errors, numerical defects
and the permitted measurement schedule. The matrix need not be $e^{Hh}$ unless
the sampled-to-continuous defect has also been bounded. A valid finite-time
enclosure is useful even without a sampled exponential-stability theorem.
Forecast bounds $\bar b_{i|k}$ require a comparison valid for \emph{every}
permitted future measurement/update history. Without that future timing
contract, propagate the true reachable set without prospective measurement
shrinkage.

Append independently valid yaw, body-velocity, gyro and absolute-position
enclosures when building the ego state set. For example, a GNSS position ball
of radius $r_p$ measured $a$ seconds ago gives radius $r_p+v_{\max}a$ about its
old measurement, not about an unrelated current estimate. Recentring requires
the displacement between those centers. The repository already publishes
current error bounds with \texttt{futurePredictionIncluded=false}. Publication
is not horizon propagation.

Sharma's exponential theorem supplies a route to \eqref{eq:comparison} only
under its assumptions and coordinate transformations. Its yaw $H_\infty$
input--output energy statement, by itself, is not a pointwise state enclosure.
A storage inequality plus an initial storage bound and a specified disturbance
energy bound could supply one; this must be derived rather than inferred from
the gain alone.

\section{Use the relative NRMM geometry, without double counting}
Let $J=\left[\begin{smallmatrix}0&-1\\1&0\end{smallmatrix}\right]$ and define
\begin{equation}
\rho=R_E^\top(p_C-p_E),\quad q=R_E^\top v_C,\quad s=R_E^\top a_C.
\end{equation}
Writing the ego body velocity as $v_E$ and yaw rate as $r_E$, the exact target
equations in the rotating ego frame are
\begin{equation}\label{eq:nrmm}
\dot\rho=q-v_E-r_EJ\rho,\qquad
\dot q=s-r_EJq,\qquad \dot s=\Phi(q,s)-r_EJs,
\end{equation}
\begin{equation}\label{eq:phi}
V_C=\norm q,\quad A_C=\frac{q^\top s}{V_C},\quad
\Omega_C=\frac{(Jq)^\top s}{V_C^2},\qquad
\Phi=-\Omega_C^2q+\frac{3A_C\Omega_C}{V_C}Jq.
\end{equation}
Here $\Omega_C$ denotes the target yaw/course rate, \emph{not} a sideslip-rate
state; $\dot\beta_C=0$. The factor three follows by differentiating
$a_C=A_Ct+\kappa_CV_C^2Jt$ with constant
$\kappa_C=\sin\beta_C/l_{rC}$. The rotating-frame terms are transport, not
independent target maneuvers.

On the positive-speed, $|\beta_C|<\pi/2$ physical domain,
\begin{equation}\label{eq:relativeYaw}
\beta_C=\arcsin\!\left(\frac{l_{rC}(Jq)^\top s}{\norm q^3}\right),\qquad
\psi_R=\arg(q)-\beta_C.
\end{equation}
The image of the joint $(q,s)$ set under \eqref{eq:relativeYaw} gives a relative
body-yaw set; interval/Lipschitz evaluation must enclose this entire image. If
$b_q<\norm{\hat q}$, the course error is at most
$\arcsin(b_q/\norm{\hat q})$; add a valid sideslip reconstruction error, with
angle wrapping. If speed can reach zero, this inverse supplies no body-heading
certificate. A saturated estimate does not establish a saturated true angle.

A fully explicit enclosure is available without a new observer state. Suppose
$\nu:=\norm{\hat q}-b_q>0$, $S_*:=\norm{\hat s}+b_s$, and truth satisfies
$|\beta_C|\le\beta_{\max}<\pi/2$. Reconstruct the center using
\begin{equation}
\hat\beta_C=\arcsin\!\left(\operatorname{clip}_{[-\sin\beta_{\max},\sin\beta_{\max}]}
\frac{l_{rC}(J\hat q)^\top\hat s}{\norm{\hat q}^3}\right).
\end{equation}
On the convex velocity ball, the unclipped argument has derivative bounds
$4l_{rC}S_*/\nu^3$ in $q$ and $l_{rC}/\nu^2$ in $s$. Clipping is
nonexpansive and the clipped arcsine has Lipschitz constant
$1/\cos\beta_{\max}$. Hence one sound bound is
\begin{equation}\label{eq:explicitYawBound}
b_\beta=\frac{l_{rC}}{\cos\beta_{\max}}
\left(\frac{4S_*b_q}{\nu^3}+\frac{b_s}{\nu^2}\right),\qquad
b_{\psi_R}=\min\!\left\{\pi,\arcsin\frac{b_q}{\norm{\hat q}}+b_\beta\right\}.
\end{equation}
The ball avoids zero throughout every connecting segment, so this argument
does not presume convexity of a positive-speed annulus. Evaluating the joint
ellipsoid image can improve this conservative component bound. An invalid
true sideslip premise cannot be made valid by clipping the estimate.

Let $B_E,B_C$ be body rectangles, including any reference-point offsets. Rigid
motion invariance gives the exact identity
\begin{equation}\label{eq:relativeDistance}
\dist(\mathcal V_E,\mathcal V_C)
=\dist\bigl(B_E,\rho+R(\psi_R)B_C\bigr).
\end{equation}
Thus the observer's relative-position and relative-yaw uncertainty can be used
directly. Common absolute translation and yaw do not enter this instantaneous
collision expression. Adding separate inertial ego and target position balls
can count this common uncertainty twice. Absolute pose uncertainty still enters
path tracking and any global-coordinate constraint; ego velocity/rate errors
still enter future \eqref{eq:nrmm}. This cancellation does not discard those
effects or assert independence of the relative errors.

\section{Information update and prediction are different operations}
\subsection{The information state and a monotone update}
Let $z$ contain the physical ego and target variables, and any constant unknown
parameters retained in prediction. Let
\begin{equation}
I_k=(\eta_k,b_k,\Z_k,u_{k-1},t_k)
\end{equation}
contain the available observer memory $\eta_k$, certificate data, a nonempty
set of possible true states, previous input and time. At initialization,
intersect the observer sets with the physical model/domain and any other
independently certified information. The essential property is
\begin{equation}\label{eq:truth}
z_k\in\Z_k.
\end{equation}
Let $\mathcal P_h(\Z,u)$ be a sound outer reachable set for the exact held-input
plant, including disturbances and numerical enclosure errors. Use
\begin{equation}\label{eq:intersection}
\Z_{k+1}=\mathcal P_h(\Z_k,u_k)\cap
\mathcal O(\eta_{k+1},b_{k+1})\cap\mathcal D_{k+1}.
\end{equation}
$\mathcal O$ is the new current observer enclosure and $\mathcal D$ contains
other valid information. An unavailable channel contributes the whole space,
not an invented finite radius; a predicted loss of target observability may
therefore leave only its propagated motion set. Each set contains truth,
hence their intersection does.
In particular
\begin{equation}\label{eq:nesting}
z_{k+1}\in\Z_{k+1}\subseteq\mathcal P_h(\Z_k,u_k).
\end{equation}
The estimate may move in either direction; the physical successor inclusion is
what matters. An empty intersection contradicts at least one premise and is
not repaired mathematically by silently increasing a safety budget.

An outer box of the intersection need not remain inside the old prediction.
For two unit disks centered at $(0,0)$ and $(1,0)$, the lens has bounding box
$[0,1]\times[-\sqrt3/2,\sqrt3/2]$; a corner of that box lies outside the first
disk. Preserve the intersection as constraints, or separately certify inclusion
of a reduced representation in the already certified tube. Simple recentering
of an unchanged physical set is harmless; unjustified enlargement is not.
Use an inclusion-preserving reach operator, or retain the old tube constraints
when a new approximation is not monotone.

\subsection{Why an observer decay rate is not a prediction decay rate}
At time $k$, the future position set of a target with uncertain constant speed
can grow even if a continuously observed estimator is exponentially stable.
For $p(0)\in[-r_p,r_p]$, $v\in[-r_v,r_v]$, the exact current-conditioned
position radius is $r_p+tr_v$, not $r_pe^{-\lambda t}$.
Future measurements can reduce \emph{future posterior} error, but their centers
and values are unknown now. There are two sound operations:
\begin{enumerate}[nosep]
\item Propagate the current true-state set under the plant with common
open-loop inputs. No future observer contraction shrinks that reachable set.
\item Predict causal output feedback and all possible measurement outcomes.
Bounds such as \eqref{eq:sampled} may bound error about each future estimate;
the family of possible estimates must also be covered.
\end{enumerate}
The second permits less conservative controlled ego tubes. It does not turn
an uncontrolled target into a contracting system. Constant $A_C,\beta_C$ are
retained along each possible target path; they are not reselected at every
prediction step. A model with bounded nonzero defect must propagate that defect
instead of claiming exact constant-parameter paths.

\subsection{A causal tube that exposes the observer contribution}
Use a causal estimate-feedback policy class. A useful finite-dimensional
member is
\begin{equation}\label{eq:policy}
u_i=v_i+K_i(\hat x_{E,i}-c_{E,i}),
\end{equation}
with preselected gains, nominal centers $c_i$ and decision variables $v_i$.
Only estimates, never unknown true states, enter this policy. These gains
describe the future response of the same controller; they are not a separate
executable backup. At the first node $c_{E,0}=\hat x_{E,k}$ is convenient, so
the feedback term is zero, but forcing this recentering is not required and
must not discard a previously feasible policy.

Linearize the joint hold map once at $(\bar z_i,\bar u_i)$, obtaining $A_i,B_i$.
With $C_Ez=x_E$, $d_i=z_i-c_i$, and estimation error
$\epsilon_{E,i}=x_{E,i}-\hat x_{E,i}$, set
\begin{align}
c_{i+1}&=f_h(\bar z_i,\bar u_i)+A_i(c_i-\bar z_i)+B_i(v_i-\bar u_i),\\
d_{i+1}&=(A_i+B_iK_iC_E)d_i-B_iK_i\epsilon_{E,i}+w_i+r_i.
\label{eq:tubeDynamics}
\end{align}
The minus sign multiplying observer error follows directly from
$\hat x_E-c_E=C_Ed-\epsilon_E$. It is not a freely chosen disturbance channel.
If $d_i\in\E_i$, $\epsilon_{E,i}\in\E^o_{E,i}$, $w_i\in\W_i$ and
$r_i\in\mathcal R_i$, a sound recursion is
\begin{equation}\label{eq:tube}
\E_{i+1}\supseteq(A_i+B_iK_iC_E)\E_i\oplus
(-B_iK_i)\E^o_{E,i}\oplus\W_i\oplus\mathcal R_i.
\end{equation}
Dropping dependence between terms is conservative, not unsound. For component
boxes with radii $r^z_i,b^E_i$, it suffices that
\begin{equation}\label{eq:boxTube}
r^z_{i+1}\ge |A_i+B_iK_iC_E|r^z_i+|B_iK_i|b^E_i+
\bar w_i+\bar r_i.
\end{equation}
Absolute value is entrywise. The current observer enclosure initializes the
tube, and future worst-case observer bounds enter its disturbance term. For
$K_i=0$, this reduces to ordinary reachable-set propagation. The closed-loop
plant tube may contract if its controlled dynamics do; the observer decay
constant is not substituted for $A_i$.

\section{Certified one-linearization collision constraints}
\subsection{Remainders over the complete decision domain}
Suppose the Jacobian of scalar map $f_{h,j}$ has Lipschitz constant $L_j$ on
the relevant convex neighborhood of the anchor. Then
\begin{equation}\label{eq:taylor}
|r_{i,j}|\le\tfrac12L_j
\norm{(z_i-\bar z_i,u_i-\bar u_i)}^2.
\end{equation}
The neighborhood covers nominal decision corrections, propagated true-state
uncertainty, output-feedback deviations and the entire connecting segment.
A bound covering only the estimation error at the anchor is insufficient.
Integration error between RK4 and the physical hold map is additional unless
RK4 itself is explicitly the entire claimed plant. A nonsmooth tire regime
requires a valid piecewise/interval enclosure, not an unjustified global
Hessian. A measured 10-mm mismatch on one rollout is not \eqref{eq:taylor}.

Freeze valid $\mathcal R_i$, tube supports and collision normals before the
convex solve, or use sound convex upper bounds depending on its variables.
This can give one conservative conic program without within-frame SCP.
Repeated linearization is not required by the safety proof; valid setwise
remainders and feasibility are. Whether sufficiently tight bounds fit a
100-ms budget is an experimental question, not a consequence of this theorem.

\subsection{A complete relative-rectangle row}
For a fixed unit normal $n$, separation in the ego frame is certified by
\begin{equation}\label{eq:plane}
\min_j n^\top\bigl(\rho+R(\psi_R)v_{C,j}\bigr)-h_{B_E}(n)\ge d_0,
\qquad h_{B_E}(n)=\max_{v\in B_E}n^\top v.
\end{equation}
Every ego point then has projection at most $h_{B_E}(n)$, and every target
point at least $d_0$ beyond it; Cauchy--Schwarz proves Euclidean distance at
least $d_0$. Normals can vary across times but are fixed during one solve.
This is a sufficient separating-plane condition, not an equivalence for an
arbitrarily selected normal. A zero multiplier from an overlapping ordinary
distance dual proves no positive clearance.

Let $p_r=(\rho,\psi_R)$, with anchor $\bar p_r$, affine center correction
$\Delta p_r$, and certified pose-error set $E_r$. For vertex $j$, define
\begin{equation}
g_j^0=n^\top(\bar\rho+R(\bar\psi_R)v_{C,j})-h_{B_E}(n),\quad
a_j^\top=[\,n^\top,\ n^\top R(\bar\psi_R)Jv_{C,j}\,].
\end{equation}
If $|\Delta\psi_R|\le t_\psi$ and $|e_{\psi_R}|\le b_\psi$, a sufficient row is
\begin{equation}\label{eq:robustRow}
\boxed{
g_j^0+a_j^\top\Delta p_r-\supp{E_r}{-a_j}
-\tfrac12\norm{v_{C,j}}(t_\psi+b_\psi)^2
-\varepsilon_i^{\rm hold}-\varepsilon_i^{\rm num}
\ \ge\ d_0-\xi_i.}
\end{equation}
The orientation remainder follows from
$|n^\top R(\psi)J^2v|\le\norm v$. Any nonlinear reconstruction of $p_r$ from
the joint state needs its own enclosure before $E_r$ is used. Pose-model errors
from \eqref{eq:tube} are already in $E_r$; do not add the same error twice.
$\varepsilon^{\rm num}$ bounds admitted numerical residuals. For ellipsoids,
\eqref{eq:support} computes the directional tightening analytically; for boxes
it is $|a_j|^\top r$. Exact maximization of rectangle supports over an angular
interval is another valid, possibly tighter, certificate.

\subsection{Whole-hold safety, not just node safety}
Let $d(t)$ be the actual rectangle distance on a hold. If the reference-point
speeds and body yaw rates are bounded throughout it, then
\begin{equation}\label{eq:holdLipschitz}
|d(t)-d(s)|\le L_d|t-s|,\quad
L_d=\bar v_E+\bar v_C+R_E\bar r_E+R_C\bar r_C,
\end{equation}
where $R_E,R_C$ bound vertex distance from the corresponding reference point.
This follows from Hausdorff motion bounds for each body and the triangle
inequality for set distance. For nodes at start, midpoint and end, every time
is within $h/4$ of a node. Thus take
\begin{equation}\label{eq:holdReserve}
\varepsilon_i^{\rm hold}=L_dh/4.
\end{equation}
More nodes or a direct interval reachable-body enclosure can reduce this
conservative reserve. The midpoint must be an enclosed physical midpoint, not
an endpoint interpolant with unbounded interpolation error. Node certificates
with $\xi_i=0$ and \eqref{eq:holdReserve} prove \eqref{eq:goal} throughout the
hold. If different slacks meet at a boundary, each adjacent hold must have its
own valid node inequalities. A positive slack permits clearance below $d_0$;
it is not a certificate of the requested margin.

\section{The terminal set: separation and the road}
Let $M$ be the last predicted node, with ego position $p_E$, yaw $\psi_E$,
body velocity $v_E$, and target position $p_C$ and velocity $w_C$. The terminal
set is
\begin{equation}\label{eq:terminal}
(p_E-p_C)^\top\bigl(R(\psi_E)v_E-w_C\bigr)\ge0,\qquad
\mathcal V_E\subseteq\text{road},
\end{equation}
together with the collision rows of every predicted node, which the endpoint
also carries. No distance from the target and no invariant cruise core are
required. The horizon length $M$ is chosen by the anchor rollout: the first node
at or after the prefix length at which the anchor separates.

\begin{proposition}[Clearance after the horizon]\label{prop:separating}
Suppose that after $t_M$ the relative velocity $v_{\rm rel}=R(\psi_E)v_E-w_C$
and both body orientations stay constant, and that \eqref{eq:terminal} and the
endpoint collision row hold at $t_M$. Then the reference-point distance is
nondecreasing for $t\ge t_M$, and with circumscribing disks of radii $R_E,R_C$
the rectangle clearance after $t_M$ is at least
$\norm{p_E-p_C}(t_M)-R_E-R_C$.
\end{proposition}
\begin{proof}
With $\Delta(t)=p_E-p_C$, $\tfrac{d}{dt}\norm\Delta^2=2\Delta^\top v_{\rm rel}$
and $\tfrac{d^2}{dt^2}\norm\Delta^2=2\norm{v_{\rm rel}}^2\ge0$. The first
derivative is nonnegative at $t_M$ and nondecreasing, so $\norm\Delta$ never
decreases. The disk bound follows from the triangle inequality.
\end{proof}
The premise is restrictive. A braking or accelerating target, a target whose
constant sideslip turns it, and the ego's own return to the path all change
$v_{\rm rel}$ after $t_M$ and can close the distance again. The next sample
then predicts the encounter again inside its own horizon; that is a per-sample
statement, not invariance. The set \eqref{eq:terminal} is therefore not
controlled invariant, and no recursive-feasibility argument rests on it. An
earlier version of this document built an information-state invariant terminal
family with a local ego core; that design is not used and has been removed.

\section{Robust PCBF and per-sample safety}
\subsection{The robust finite-horizon problem}
A prediction policy $\Pi=(\pi_0,\ldots,\pi_{M-1})$ is causal: $\pi_i$ uses
information available at that future node. Let $\mathfrak I_i^\Pi(I)$ be all
possible future information states generated from $I$. Use common nonnegative
prefix slacks $\xi_0,\ldots,\xi_{N-1}$; extend $\xi_i=0$ for $N\le i<M$.
The feasible witness set $\mathcal F_M(I)$ consists of policies and slacks
satisfying, for every information branch and every state/disturbance in it,
\begin{align}\label{eq:robustProblem}
&\text{dynamics, observer/update contract, and hard physical constraints},\nonumber\\
&\dist(\mathcal V_E(t),\mathcal V_C(t))\ge d_0-\xi_i
\quad(t\in[t_{k+i},t_{k+i+1}]),\\
&\text{the terminal set \eqref{eq:terminal} at node }M.\nonumber
\end{align}
Any additional softened constraints use specified scales and the same stage
slacks; this document uses distance units for collision slack. Softening a
model-evaluability domain is not allowed.
Define
\begin{equation}\label{eq:pcbf}
V_B^*(I)=\inf_{(\Pi,\Xi)\in\mathcal F_M(I)}\sum_{i=0}^{N-1}\xi_i,
\quad S(\Xi)=\sum_{i=0}^{N-1}\xi_i.
\end{equation}
$V_B^*$ is defined on information, not just the estimated vehicle state.

\begin{proposition}[Per-sample robust safety]\label{prop:sampleSafety}
Under Assumption~\ref{ass:contract}, suppose that at sample $k$ the problem
\eqref{eq:robustProblem} has a witness with $\xi_{0|k}=0$, every constraint
certificate is sound with its residuals enclosed, and its first input is
applied. Then the applied hold satisfies \eqref{eq:goal}. If the witness has
$S=0$, every predicted hold of every admitted branch satisfies \eqref{eq:goal},
and Proposition~\ref{prop:separating} applies after $t_{k+M}$ under its premise.
\end{proposition}
\begin{proof}
Truth containment follows from Proposition~\ref{prop:observer} and
\eqref{eq:intersection}, then from sound reachable-set propagation. The
realized first hold is an admitted branch, and its whole-hold robust
constraints with $\xi_0=0$ imply \eqref{eq:goal}.
\end{proof}
Nothing here guarantees that sample $k+1$ again has a zero-slack witness. A
shifted witness need not exist, because \eqref{eq:terminal} is not invariant.
Each sample therefore minimizes prefix slack again; a positive slack witness is
a restoration object, not a proof of clearance $d_0$, and its first hold only
has the certified lower bound $d_0-\xi_0$. If a numerical solver fails to return
a command within a deadline, the controller reports that no solution exists; a
runtime completion premise must be stated separately.

\subsection{What the implementation realizes (October 5, 2026)}
The executable controller does not enumerate information branches. Every hold
$i$ of its plan starts from the \emph{current} published enclosure $G_0$ and
propagates it through that hold only, $G(t_i+\tau)=\Phi_i(\tau)G_0$ for
$\tau\in[0,h]$, with the common-pose cancellation taken about the ego at $t_i$.
The target set restarts at $t_i$ from the target's predicted state, for the
elapsed time $\tau$. For $i=0$ this is the propagation that
Proposition~\ref{prop:sampleSafety} needs for the applied hold, sampled at the
hold start, midpoint and end (subject to the remainder and whole-hold premises
of Section~4). For $i\ge1$ it is an information-state approximation. It
assumes that the estimator's enclosure at $t_i$ equals the present one, and
that the estimate at $t_i$ lies on the plan. Neither the family of possible
future estimates nor the target's forecast error beyond one hold is covered,
so the $S=0$ statement about every predicted hold is not claimed. The
open-loop image $G_{i+1}=A_iG_i$ over the variable horizon (up to 512 holds)
was sound for an open-loop input sequence but made the hard road and state
rows infeasible in noisy encounters.

The hard physical constraints of \eqref{eq:robustProblem} include, at both
ends of every hold under that hold's input $b$, the estimator's sideslip cone
$|v_y|\le\tan(b_\beta)v_x$, robust to the enclosure. The rear-adhesion cone
$\norm{[(v_y-l_rr)/k;\ \bar v_xb]}\le v_x$, with $k=3\mu_rF_{z,r}/C_r$, keeps
the rear Fiala tire below saturation with capacity $\sqrt{1-b^2}$. It is
imposed on the estimate as a handling and observability condition. It is not
part of the safety certificate. The
terminal set \eqref{eq:terminal} adds road recovery: with $w_e$ the velocity
toward road edge $e$ and $d_e$ the outermost corner's distance to it,
$\max(w_e,0)^2\le2a\,d_e$ for a declared lateral deceleration $a$. Inside the road
alone admits endpoints that leave it shortly after the horizon. This is still
not an invariance argument for the road.

\section{One CLF: nominal dissipation with an explicit uncertainty supply}
\emph{Implementation note (October 5, 2026).} The executable CLF row is
$V(\hat x^+)\le\rho V(\hat x)+s$ on the estimate. It does not use the supply
derived below: with the published enclosures, the robust budget is zero near
the path while the successor's enclosure is not. The worst-case row then
cannot be met, and its slack minimization made the commands chatter. Recovery
of the true state is therefore only practical, to a neighborhood set by the
estimation error.

\subsection{The same function during avoidance and recovery}
Let
\begin{equation}\label{eq:clf}
e(x_E)=\begin{bmatrix}e_y&e_\psi&v_x-v_{x,\rm ref}&
v_y-v_{y,\rm ref}&r-r_{\rm ref}\end{bmatrix}^{\!\top},\quad
V=e^\top P e,\quad\ell=e^\top Qe,\quad P,Q\succ0.
\end{equation}
This is the existing transverse function, independent of the target. A smooth
path chart and a verified recovery domain are required. Positive definiteness
of a matrix synthesized at one trim does not prove a global nonlinear CLF. If $P$
depends on path phase, all successor expressions include the successor $P$;
freezing it without a bound is invalid. Below $P$ is fixed on the verified domain.

The desired robust first-hold row is
\begin{equation}\label{eq:robustClf}
\sup_{z\in\Z_k,\,w\in\W_k}
\{V(x_E^+)-V(x_E)+\eta\ell(x_E)\}
\le\delta_k+\sigma_k,\qquad \delta_k\ge0,\quad\eta>0.
\end{equation}
Here $\sigma_k$ is a \emph{derived} disturbance supply, not an optimized slack,
and $\delta_k$ is the single CLF slack. The supremum preserves the dependence
between a true current state and its successor. Safety constraints keep
priority through a cap on the prefix slack sum at its primary optimum. No target-range switch changes $V$.

\subsection{Deriving the supply, rather than choosing a convenient tolerance}
Suppose a common estimate-feedback action on the recovery domain has
\begin{equation}\label{eq:nominalContraction}
\sqrt{V(F_E(x_E,\kappa(x_E)))}\le a\sqrt{V(x_E)},\qquad 0\le a<1,
\end{equation}
and the effect of using $\kappa(\hat x_E)$, plant disturbance and hold-model
defect is bounded by $D_k$ in the successor $P$ norm. For example a valid
uniform bound is $D_k=L_{F,u}L_\kappa b_{E,k}+d_{E,k}$, in consistent norms.
For any $a^2<\lambda<1$,
\begin{equation}\label{eq:young}
(a\sqrt V+D_k)^2\le\lambda V+
\frac{\lambda}{\lambda-a^2}D_k^2.
\end{equation}
Indeed the difference is
$(\lambda-a^2)(\sqrt V-aD_k/(\lambda-a^2))^2$.
Choose $\eta$ so that $\eta Q\preceq(1-\lambda)P$ and set
\begin{equation}\label{eq:supply}
\sigma_k=\frac{\lambda}{\lambda-a^2}D_k^2.
\end{equation}
Then the common action satisfies \eqref{eq:robustClf} with $\delta_k=0$.
This does not show it is collision-safe, terminal-feasible or inside a
particular RTI decision domain. Zero CLF slack is guaranteed admissible only
where those requirements are jointly compatible.

\subsection{A tractable conservative SOC row}
Let $P_-=P-\eta Q\succeq0$. Precompute a sound lower bound
\begin{equation}\label{eq:clfLower}
0\le\underline W_k\le\inf_{z\in\Z_k}e(x_E)^\top P_-e(x_E).
\end{equation}
For an enclosed current error $e=\hat e+d_e$,
one choice is $(\max\{0,\norm{P_-^{1/2}\hat e}-r_{-,k}\})^2$ with
$r_{-,k}\ge\sup\norm{P_-^{1/2}d_e}$. Include nonlinear path-projection
remainders in $d_e$. Let $e_1^c$ be the affine next-error center and
$r_{P,1}$ a valid bound on $\norm{P^{1/2}(e_1-e_1^c)}$ over the whole
uncertainty and decision domain. Introduce $t\ge0$ and impose
\begin{equation}\label{eq:clfSoc}
\norm{P^{1/2}e_1^c}\le t,\qquad
(t+r_{P,1})^2\le\underline W_k+\sigma_k+\delta_k.
\end{equation}
These are ordinary/rotated SOC constraints when the radii are fixed. The
triangle inequality and \eqref{eq:clfLower} imply \eqref{eq:robustClf}.
This sufficient bound can be more conservative than the joint supremum;
therefore physical feasibility of zero slack does not ensure its feasibility
in \eqref{eq:clfSoc}. A joint quadratic robust bound can be tighter but may
require an SDP or a different certified approximation. The computation is not
automatically a QP with only affine constraints.

\subsection{A proximal tie that does not create a nonzero recovery floor}
The current controller adds a bounded input-increment term relative to the
linearization input, not zero. A fixed additive weight can buy positive CLF
slack even when zero is feasible. To retain one solve and remove that
mathematical obstruction, let $0\le R(\Delta U)\le1$ on the decision domain,
and compute
\begin{equation}\label{eq:ellLower}
0\le\underline\ell_k\le\inf_{z\in\Z_k}\ell(x_E).
\end{equation}
Use the single secondary objective
\begin{equation}\label{eq:objective}
\boxed{\min\ \delta_k+\epsilon\,\underline\ell_k R(\Delta U),
\qquad 0\le\epsilon<\eta,}
\end{equation}
under the robust safety budget and all constraints above. Scaling this entire
objective by a positive known constant is harmless. Scaling just one term
changes the argument.

\begin{proposition}[One-solve CLF recovery with a vanishing tie]\label{prop:clf}
Suppose from some sample onward a zero-CLF-slack feasible candidate exists in
the actual optimization, its optimum is attained, and its objective is solved
with gap at most $\tau_k\ge0$. An issued optimizer satisfying
\eqref{eq:robustClf} and \eqref{eq:objective} obeys
\begin{equation}\label{eq:recovery}
V_{k+1}-V_k\le-(\eta-\epsilon)\ell_k+\sigma_k+\tau_k.
\end{equation}
If $\ell\ge\mu V$ on a forward invariant bounded recovery domain, with
$0<c:=(\eta-\epsilon)\mu\le1$, then
\begin{equation}\label{eq:iss}
V_{k+j}\le(1-c)^jV_k+
\sum_{s=0}^{j-1}(1-c)^{j-1-s}(\sigma_{k+s}+\tau_{k+s}).
\end{equation}
In particular $\limsup V_k\le(\bar\sigma+\bar\tau)/c$ for uniformly bounded
supplies, and $V_k\to0$ if the supplies tend to zero.
\end{proposition}
\begin{proof}
A zero-slack feasible point has objective at most
$\epsilon\underline\ell_k$. Nonnegativity of $R$ gives
$\delta_k\le\epsilon\underline\ell_k+\tau_k\le\epsilon\ell_k+\tau_k$
for the actual state. Insert this in \eqref{eq:robustClf} and iterate the scalar
contraction inequality. A stable scalar convolution of a bounded sequence
tending to zero also tends to zero, proving the last claim.
\end{proof}
The objective minimizes CLF relaxation under the safety budget. It is not
an assertion of minimum Euclidean distance to a nominal input. The proximal
term remains centered on the linearization input, not on zero, and introduces
neither a second CLF nor a separate path-deviation penalty.
Thus exact zero numerical slack is not necessary for genuine dissipation:
the allowable optimization tie vanishes with the available tracking error.
For an exact lexicographic minimum of $\delta$, simply set $\epsilon=0$ or
use a solver that enforces that hierarchy. A fixed positive solver tolerance,
persistent model/noise bounds, or an unvanishing proximal tie generally yields
a practical neighborhood rather than exact nominal convergence.

The nominal manifold has arbitrary longitudinal phase, so $V\to0$ means
return to path-following cruise, not arrival at a scheduled position. Nothing
in the proof imposes a fixed recovery duration. Conversely, collision threat
disappearing alone does not prove membership in a global recovery domain or
zero slack in a conservative convex inner problem. Those conditions must be
certified or assessed separately.

\section{The resulting online problem and the precise implementation boundary}
At each sample perform the following operations, all within the same controller:
\begin{enumerate}[leftmargin=*,itemsep=3pt]
\item Obtain timestamped observer sets and intersect them with the certified
previous successor set using \eqref{eq:intersection}.
\item Shift the previous plan as the search anchor; without one, generate an
APF search seed. Neither has independent safety permission.
\item Construct one affine center model and valid setwise remainders/tube
supports, preserving an inherited feasible certificate when it exists.
\item Minimize prefix safety slack at every sample, then solve the \emph{same}
CLF secondary problem under a cap at that optimum; zero PCBF slack does not
skip CLF optimization.
\item Solve \eqref{eq:objective} subject to affine center dynamics,
\eqref{eq:tube} or \eqref{eq:boxTube}, \eqref{eq:robustRow}, hard tail, the
terminal rows \eqref{eq:terminal}, and \eqref{eq:clfSoc}. Apply its first
estimate-based input.
\end{enumerate}
With fixed gains, directions and remainder radii, and the terminal rows
linearized at the anchor, the center problem is a convex conic QP/SOCP.
General information-state reachability is not claimed to be an SOCP. Robust overapproximation quality, certificate closure and
first feasibility remain substantive design requirements. A trusted reference
seed can improve conditioning and direction choice, but does not replace any
uncertainty or remainder term.

\begin{center}
\fbox{\begin{minipage}{0.94\linewidth}
\textbf{Completed mathematical connection.}
Observer comparison $\Rightarrow$ current truth-containing sets
$\Rightarrow$ causal reachable tubes and nested information updates
$\Rightarrow$ whole-hold robust rectangle constraints
$\Rightarrow$ a per-sample robust predictive safety witness.
A zero-slack witness gives $d_0$ clearance on the applied hold; no recursive
feasibility is claimed.
The one CLF and its disturbance supply give practical recovery under persistent
uncertainty, and asymptotic nominal recovery when uncertainty/numerical supply
vanishes and zero CLF slack remains feasible in the recovery domain.
\end{minipage}}
\end{center}

\subsection{Experimental implementation and remaining certification gap}
The current implementation tests an open-loop approximation of the ideal
construction above. It does not enforce all assumptions of
Proposition~\ref{prop:sampleSafety}. In particular, a newly published target acceleration or
sideslip estimate may change at every sample; each prediction freezes those
new estimates, without requiring a nested information update.

\begin{longtable}{p{0.27\linewidth}p{0.66\linewidth}}
\toprule
Implemented component & Behavior and qualification \\\midrule
\endhead
Observer publication
& A joint current ego generator retains yaw/body-velocity correlation from
GNSS. Measurement-history enclosures tighten target parameter intervals.
Their validity still depends on the declared observer and sensor premises. \\
Prediction geometry
& The common initial rigid pose is retained through propagation and cancelled
only in relative collision geometry. Directional target supports propagate
constant-parameter uncertainty without assuming future target measurements. \\
Ego uncertainty tube
& The current correlated generator is propagated open loop through the affine
hold maps, $G_{i+1}=A_iG_i$. No ancillary feedback is assumed, so the issued
inputs carry no uncertainty of their own. This is an experimental envelope,
not a proved future bound. \\
Terminal set
& The implemented terminal set is separation from the target estimate, with
the endpoint also under the collision rows, and the ego rectangle inside the
road, with ego generator supports subtracted. No invariance is claimed. \\
Safety and one CLF
& With nonzero uncertainty, mean collision rows are hard and uncertainty
tightenings may use slack throughout the finite horizon. The same quadratic
CLF has uncertainty padding and a tracking-error-scaled proximal tie. \\
RTI admission
& One model is built per shifted or fresh APF anchor (a shifted problem infeasible inside the trust region is re-solved once from a
fresh rollout); the PCBF stage is solved in every frame, and a frame without a solution is reported. The completed
solution is issued without a nonlinear replay. The next plan innovation is split exactly
into an observation part (posterior minus a nonlinear replay from the plan's own
prediction) and a model part (replay minus the affine prediction); only the
model part sets the input trust scale. This
does not certify the nonlinear plan or continuous-time robust separation. \\
\bottomrule
\end{longtable}

Writing $d_i=x_i-\bar x_i$, the implemented propagation is the open-loop
image $d_{i+1}=A_id_i$ from the current enclosure $d_0=\varepsilon_0$. The causal
policy class of \eqref{eq:policy} is not used: with $K_i=0$ the tube reduces to
reachable-set propagation, which is sound but more conservative than a
feedback-aware tube. Validated nonlinear remainders would still be required to
promote this propagation to \eqref{eq:tube}.

Every sample minimizes the prefix slack and then the CLF slack under a cap at
that optimum. The achieved value is not $V_B^*$, and changing observer sets do
not establish cross-sample monotonicity. Positive uncertainty slack explicitly
forfeits a hard robust collision certificate; even hard mean rows concern the
sampled affine model.

Thus the implementation is an empirical output-feedback controller with
uncertainty-aware planning, not an implementation of every premise of the
theorem. The finite identity checks in
\texttt{verifyObserverRobustConnection.m} do not establish global
Fiala remainder bounds or a runtime deadline. Collision and recovery experiments must retain that scope.

\begin{thebibliography}{9}
\bibitem{Sharma} G. Sharma, H. Alai, and R. Rajamani,
``Simultaneous ego-vehicle state estimation and vehicle trajectory tracking
using a multistage high gain observer,'' \emph{Transportation Research Part C},
vol. 182, 105411, 2026.
\href{https://doi.org/10.1016/j.trc.2025.105411}{doi:10.1016/j.trc.2025.105411}.
Local reference PDF, Assumption 1 and Theorems 1--2, PDF pp. 6--8.
\bibitem{Huang} J. Huang, H. Wang, K. Margellos, and P. Goulart,
``Predictive Control Barrier Functions: Bridging model predictive control
and control barrier functions,'' \emph{European Control Conference},
pp. 2623--2629, 2025.
\href{https://doi.org/10.23919/ECC65951.2025.11187291}{doi:10.23919/ECC65951.2025.11187291};
\href{https://arxiv.org/abs/2502.08400}{arXiv:2502.08400}, Sec. III, Eq. (8).
\bibitem{Kohler} J. K\"ohler, M. A. M\"uller, and F. Allg\"ower,
``Robust output feedback model predictive control using online estimation
bounds,'' 2021 preprint.
\href{https://arxiv.org/abs/2105.03427}{arXiv:2105.03427}, Secs. III--IV.
\bibitem{Local} Collision Avoidance repository,
\emph{Observer ISS Theory}, \texttt{estimator/OBSERVER\_ISS\_THEORY.md};
\emph{One Analytic Nominal CLF}, \texttt{controller/NOMINAL\_CLF.md}; and
\emph{Observer-to-Robust-PCBF Gap},
\texttt{report/OBSERVER\_ROBUST\_PCBF\_GAP\_20261003.tex}.
Working-source hashes and finite-check provenance accompany this derivation.
\end{thebibliography}
\end{document}
